\section{Implement Trie II (Count and Erase)}
\label{sec:trie-implement-2}

\problemstatement{Implement a trie supporting $\textit{insert}(word)$,
$\textit{countWordsEqualTo}(word)$ (how many times $word$ has been
inserted, counting duplicates), $\textit{countWordsStartingWith}(prefix)$
(how many inserted words, counting duplicates, begin with $prefix$), and
$\textit{erase}(word)$ (remove one occurrence of a previously inserted
$word$).}

\begin{patternbox}
The keyword shift from Section~\ref{sec:trie-implement} is from a
yes/no question (\emph{is this word present}) to a \textbf{counting}
question (\emph{how many times}, including duplicates and including
prefix counts), plus \textbf{removal}. The fix is to annotate every node,
not just mark a single boolean: track how many inserted words
\textbf{pass through} this node (as part of some longer word, or ending
exactly here), and how many \textbf{end exactly} at this node.
\end{patternbox}

\subsection*{Understanding the problem carefully}

Every node along the path of an inserted word is, by definition, a
prefix of that word -- so if we increment a $\textit{countPrefix}$
counter at \emph{every} node visited during an insertion, that counter at
any given node ends up holding exactly the number of inserted words
(counting duplicates) that share this far as a common prefix. Separately,
incrementing an $\textit{countEnd}$ counter only at the \emph{final} node
of each inserted word tracks exactly how many times that complete word
itself was inserted.

\begin{ideabox}[Two Counters per Node: Prefix-Passes and Exact-Ends]
$\textit{insert}(word)$: walk the word's path, creating nodes as needed
(as in Section~\ref{sec:trie-implement}), and increment
$\textit{countPrefix}$ at \emph{every} node visited along the way
(including the final one); increment $\textit{countEnd}$ only at the
final node. $\textit{countWordsEqualTo}(word)$: walk the path; if it
exists in full, return the final node's $\textit{countEnd}$ (else $0$).
$\textit{countWordsStartingWith}(prefix)$: walk the path; if it exists in
full, return the final node's $\textit{countPrefix}$ (else $0$).
$\textit{erase}(word)$: walk the path (assuming $word$ was indeed
previously inserted), \textbf{decrementing} $\textit{countPrefix}$ at
every node visited, and decrementing $\textit{countEnd}$ at the final
node.
\end{ideabox}

\subsection*{Why two separate counters are needed, not just one}

A single counter cannot distinguish ``$\textit{prefix}$ is itself a
complete inserted word $k$ times'' from ``$\textit{prefix}$ is merely a
prefix of $k$ longer inserted words.'' For example, after inserting
``app'' once and ``apple'' once, the node at the end of ``app'' has
$\textit{countPrefix}=2$ (both words pass through it) but
$\textit{countEnd}=1$ (only ``app'' itself ends there) --
$\textit{countWordsEqualTo}(\texttt{"app"})$ must report $1$, while
$\textit{countWordsStartingWith}(\texttt{"app"})$ must report $2$, and
no single counter can serve both queries correctly.

\subsection*{Algorithm}

\begin{enumerate}
  \item \textsc{Insert}($word$): walk/create the path; increment
        $\textit{countPrefix}$ at every visited node; increment
        $\textit{countEnd}$ at the final node.
  \item \textsc{CountWordsEqualTo}($word$): walk the path; return the
        final node's $\textit{countEnd}$, or $0$ if the path does not
        fully exist.
  \item \textsc{CountWordsStartingWith}($prefix$): walk the path; return
        the final node's $\textit{countPrefix}$, or $0$ if the path does
        not fully exist.
  \item \textsc{Erase}($word$): walk the path; decrement
        $\textit{countPrefix}$ at every visited node; decrement
        $\textit{countEnd}$ at the final node.
\end{enumerate}

\subsection*{Dry run}

\dryrun{Insert \texttt{"apple"} twice, insert \texttt{"app"} once. Query
$\textit{countWordsEqualTo}(\texttt{"apple"})$,
$\textit{countWordsStartingWith}(\texttt{"app"})$. Then erase
\texttt{"apple"} once and repeat both queries.}

\begin{itemize}
  \item After both insertions: node at \texttt{"app"} has
        $\textit{countPrefix}=3$ (all three insertions pass through
        it), $\textit{countEnd}=1$ (only the single ``app'' insertion
        ends there). Node at \texttt{"apple"} has
        $\textit{countPrefix}=2$, $\textit{countEnd}=2$.
  \item $\textit{countWordsEqualTo}(\texttt{"apple"}) = 2$.
        $\textit{countWordsStartingWith}(\texttt{"app"}) = 3$.
  \item Erase \texttt{"apple"} once: decrement $\textit{countPrefix}$
        along the whole path (node at \texttt{"app"}: $3\to2$; node at
        \texttt{"apple"}: $2\to1$), and $\textit{countEnd}$ at the final
        node ($2\to1$).
  \item Now $\textit{countWordsEqualTo}(\texttt{"apple"}) = 1$.
        $\textit{countWordsStartingWith}(\texttt{"app"}) = 2$.
\end{itemize}

\subsection*{C++ implementation}

\begin{lstlisting}
#include <bits/stdc++.h>
using namespace std;

struct TrieNode {
    TrieNode* children[26] = {nullptr};
    int countPrefix = 0;
    int countEnd = 0;
};

class TrieII {
    TrieNode* root;

public:
    TrieII() { root = new TrieNode(); }

    void insert(string word) {
        TrieNode* node = root;
        for (char c : word) {
            int idx = c - 'a';
            if (!node->children[idx]) node->children[idx] = new TrieNode();
            node = node->children[idx];
            node->countPrefix++;
        }
        node->countEnd++;
    }

    int countWordsEqualTo(string word) {
        TrieNode* node = root;
        for (char c : word) {
            int idx = c - 'a';
            if (!node->children[idx]) return 0;
            node = node->children[idx];
        }
        return node->countEnd;
    }

    int countWordsStartingWith(string prefix) {
        TrieNode* node = root;
        for (char c : prefix) {
            int idx = c - 'a';
            if (!node->children[idx]) return 0;
            node = node->children[idx];
        }
        return node->countPrefix;
    }

    void erase(string word) {
        TrieNode* node = root;
        for (char c : word) {
            int idx = c - 'a';
            if (!node->children[idx]) return; // not present
            node = node->children[idx];
            node->countPrefix--;
        }
        node->countEnd--;
    }
};

int main() {
    TrieII trie;
    trie.insert("apple");
    trie.insert("apple");
    trie.insert("app");

    cout << trie.countWordsEqualTo("apple") << endl;        // 2
    cout << trie.countWordsStartingWith("app") << endl;      // 3

    trie.erase("apple");
    cout << trie.countWordsEqualTo("apple") << endl;        // 1
    cout << trie.countWordsStartingWith("app") << endl;      // 2
    return 0;
}
\end{lstlisting}

\begin{complexitybox}
\textbf{Time Complexity.} $O(L)$ for every operation.
\textbf{Space Complexity.} $O(\Sigma \cdot N)$, as in
Section~\ref{sec:trie-implement}, with two integer counters added per
node.
\end{complexitybox}

\begin{takeawaybox}
Whenever a trie problem needs counts rather than a single boolean,
annotate every node with the counters the queries actually need --
typically ``how many words pass through here'' and/or ``how many words
end exactly here'' -- and update them incrementally on insert and erase.
This is a direct generalization of the single $\textit{isEnd}$ flag from
Section~\ref{sec:trie-implement}.
\end{takeawaybox}
